\documentclass[12pt,a4paper,oneside]{article}
\usepackage[brazil]{babel}
\usepackage[utf8]{inputenc}
\usepackage{olymp}

\setcounter{problem}{2}

\begin{document}

\begin{problem}{Guerra dos Tronos}{guerra}{2}
 
Capivarus e Canideus travaram uma guerra durante 50 anos. O motivo da guerra era o tamanho do território de cada reino. Pelo bem da população dos dois reinos, os governantes de seus tronos resolveram fazer um tratado para finalizar a guerra. O tratado consiste em fazer um divisão justa, e certamente contínua, do território. Eles resolveram pedir sua ajuda para calcular o ponto de divisão do território. Depois de tantos anos de guerra, os reinos não podem lhe pagar uma viagem para ver previamente o território que será dividido. Ao invés disso, eles prepararam uma lista $a_1, a_2, \dots, a_N$ de inteiros que indicam o tamanho de cada seção do território. A seção $a_1$ é vizinha da seção $a_2$ que por sua vez é vizinha da seção $a_3$; e assim por diante. Os reinos querem uma divisão em uma seção $k$ de tal forma que $a_1 + a_2 + \dots + a_k = a_{k+1} + a_{k+2} + \dots + a_N$.

%%% Acrescentado, não havia no texto da OBI
%%%-----------
Como exemplo, considere o território da figura abaixo, que contém 4 seções, de tamanhos 5, 3, 2 e 10, nesta ordem.

\vspace{-21pt}
\begin{center}
\includegraphics[scale=0.3]{images/exercise_80_0.png}
\end{center}
\vspace{-21pt}

A solução é dividir o território após a seção 3, conforme figura abaixo. Assim, Capivarus fica com 10 de território e Canideus também.

\vspace{-21pt}
\begin{center}
\includegraphics[scale=0.3]{images/exercise_80_1.png}
\end{center}
\vspace{-21pt}

%%%------------

Sua tarefa é, dada uma lista de inteiros positivos $a_1, a_2, \dots, a_N$, determinar a seção $k$ tal que a soma dos comprimentos das seções $a_1$ até $a_k$ seja igual à soma dos comprimentos das seções $a_{k+1}$ até $a_N$.


%\Notes
%
%Observações, se tiver.

\InputFile

A entrada é composta por duas linhas. A primeira contém um valor inteiro $N$, que indica o número de seções do território. A segunda linha contém $N$ valores inteiros  $a_1, a_2, \dots, a_N$ indicando os comprimentos das seções. Restrições: $1 < N \leq 10^5$; $1 \leq a_i \leq 100, \forall i=1\dots N$.

\OutputFile

Seu programa deve imprimir uma única linha contendo um inteiro que indica a seção do território onde acontecerá a divisão.

\Notes
É garantido que sempre existe uma divisão que satisfaz as condições dos reinos.

\Example

\begin{example}
\exmp{
4
5 3 2 10
}{
3
}%
\end{example}

\begin{example}
\exmp{
9
2 8 2 8 4 4 4 4 4
}{
4
}%
\end{example}

\small \textit{Obs.: baseada na questão ``Guerra por território'' da OBI 2012} 

%Comentários sobre o exemplo, se necessário.

\end{problem}

\end{document}
